\documentclass[a4paper,10pt]{book}
\usepackage[utf8x]{inputenc}
\usepackage{graphicx}
\DeclareGraphicsRule{.eps}{eps}{}{}
\usepackage{hyperref}

\title{Controller Area Network}
\author{Federico Vaga}
\date{Giugno 2012}
\makeindex
\begin{document}
\maketitle
% stampa la licenza CC BY-NC-SA_2.5_IT
% occorre che i documento originale usi il pacchetto hyperref

\begin{center}\textbf{{\large Licenza - License}}\end{center}

\begin{figure}[hcbp]  
  \centering
  \includegraphics{licenza/CreativeCommons/BY-NC-SA_2.5_IT/cc.png}
\end{figure}

\begin{center}\textbf{Italiano}\end{center}
“Questo lavoro è licenziato sotto licenza Creative Commons - Non commerciale - Condividi allo stesso modo 2.5 Italia (CC BY-NC-SA 2.5). Per vedere una copia di questa licenza visita il sito \url{http://creativecommons.org/licenses/by-nc-sa/2.5/it/}; oppure invia una lettera a Creative Commons, 171 2nd Street, Suite 300, San Francisco, California, 94105, USA.”\newline
\begin{center}\textbf{English}\end{center}
“This work is licensed under the Creative Commons Attribuzione - Non commerciale - Condividi allo stesso modo 2.5 Italia (CC BY-NC-SA 2.5) License. To view a copy of this license, visit \url{http://creativecommons.org/licenses/by-nc-sa/2.5/it/}; or, (b) send a letter to Creative Commons, 171 2nd Street, Suite 300, San Francisco, California, 94105, USA.” \newline



\chapter{Introduzione}
\section{Storia del protocollo}
Il \textit{Controller Area Network}, abbreviato in \textit{CAN}, è un
protocollo di comunicazione seriale ideato da Bosch negli anni '80.
Questo protocollo è stato espressamente progettato per resistere a
forti interferenze dovute ad onde elettromagnetiche esterne; questo
suo punto di forza fa si che il protocollo CAN trovi largo utilizzo
nel settore automotive e nel settore dell'automazione industriale.
Il lavoro a questo protocollo iniziò nel 1983 e fu ufficialmente
rilasciato nel 1986; nel 1987 fu commercializzato il primo controller
CAN prodotto da Intel e Philips; nel 1991 Bosch rilascia le
specifiche per il CAN 2.0.
\newline

Dal 1996 negli Stati Uniti d'America tutti i veicoli commercializzati
devono sottostare allo standard \textit{OBD-II} per la diagnostica di
bordo, e il protocollo CAN è uno dei cinque protocolli consentiti. In
Europa dal 2001 è obbligo rispettare lo standard EOBD per i veicoli a
benzina, mentre dal 2004 per quelli a diesel; anche per lo standard
europeo il protocollo CAN è uno di quelli ammessi.
\newline

In protocollo CAN è riconosciuto come standard con il rilascio dell'
ISO 11898.

\section{Caratteristiche principali}
Come accennato in precedenza, il protocollo CAN è stato progettato
per essere molto robusto alle interferenze ed essere immune da errori
di trasmissione. Questa grande resistenza è dovuta all'applicazione
della specifica RS-485 del livello fisico del modello OSI. La
specifica definisce la trasmissione di un dato mediante la differenza
di potenziale su due differenti fili; quindi avremo uno stato logico
alto quando la differenza di potenziale è superiore a $0.2V$, basso
quando inferiore. La resistenza al rumore può essere aumentata
realizzando un doppino intrecciato (twisted pair).
\newline

Il protocollo garantisce una velocità di trasmissione massima di
$1\frac{Mb}{s}$ per distanze non superiori a $40m$; distanze
maggiori sono raggiungibili a velocità inferiori.
\newline

I protocollo è basato su una comunicazione broadcast, quindi tutti i
nodi ascoltano le comunicazioni che avvengono sul bus anche quando
non sono dirette a loro. In questo modo è molto semplice rimuovere o
aggiungere nodi alla rete senza dover effettuare alcuna modifica
hardware o software. La trasmissione è message-ortiended, ovvero
basata sui messaggi; ogni messaggio contiene un identificativo
univoco che identifica l'origine del messaggio all'interno della rete

\chapter{La specifica}
La standard CAN descrive i livelli fisico e data link del modello
\textit{ISO/OSI} per la connessione fra calcolatori.
\section{Livello fisico}
Il livello fisico è costituito da due cavi denominati \textit{CAN\_H}
e \textit{CAN\_L}. I livelli logici trasmessi sono definiti
\textit{dominante} e \textit{recessivo}; non c'è una corrispondenza
diretta con i classici stati logici 1 e 0, perché questo dipende da
come i nodi si interfacciano con il bus. Durante la trasmissione i
bit dominanti hanno sempre la precedenza sui bit recessivi; se
durante una trasmissione i nodi inviano un bit recessivo ma un
diverso nodo invia un bit dominante, allora il bit che effettivamente
è trasmesso sul bus è quello dominante. Tutti i nodi che
trasmettevano un bit recessivo si accorgono della presenza del bit
dominante, quindi smettono di trasmettere e rimangono in ascolto.

I nodi possono essere \textit{OR-Wired} oppure \textit{AND-Wired};
nel primo caso avremo la corrispondenza dominante:1, recessivo:0; nel
secondo caso la corrispondenza è dominante:0 recessivo:1

%FIXME figura dei circuit per fare or wired e and wired

Questo livello è diviso in tre sottolivelli: PLS, PMA, MDI
\subsection{Sottolivello PLS}
La sigla PLS corrisponde a \textit{Physical Signaling}. Questo
sottolivello si occupa della sincronizzazione e della temporizzazione
dei segnali attraverso la tecnica di \textit{bit stuffing}, ovvero
l'inserimento di un bit di controllo all'interno della comunicazione.
Nel nostro caso, ogni 5 bit viene inserito un bit di polarità
inversa; questo facilita la sincronizzazione fra i nodi.
\newline

Il \textit{bit-time}, ovvero la durata di un bit, è suddivisa in
quattro parti dette \textit{segment}, ciascuna multipla del
\textit{time-quantum} \textit{TQ}. Il TQ è l'unità di tempo usata dal
protocollo per suddividere il tempo; esso è definito mediante l'uso di
un oscillatore locale. Le quattro parti che compongono la trasmissione
di un bit sono:
\begin{description}
 \item[Syncronization] dura 1TQ ed è utilizzata per inserire fronti
di commutazione e sincronizzazione fra i nodi.
 \item[Propagation time] può durare da 1 a 8 TQ, è utilizzato per
compensare i ritardi di trasmissione/ricezione.
 \item[Phase 1 e 2] entrambe i segmenti hanno la stessa durata che
può variare da 1 a 8 TQ. Fra i due segmenti viene posto il
\textit{sample-point}, ovvero il punto in cui viene letto il livello
presente sul bus. L'ampiezza di questi due segmenti è regolata di modo
da compensare le fasi. Il sistema verifica la sincronizzazione fra
nodi durante la transizione fra recessivo e dominante e calcola la
distanza del fronte con il sample-point; questa distanza
\underline{deve} essere sempre la stessa. I fronti sono rilevati ad
ogni TQ e viene confrontato lo stato della linea con lo stato
dell'ultimo sample-point. Per mantenere costante la distanza fra il
fronte e il sample-point, i segmenti di fase devono essere
opportunamente dimensionati; è possibile definire un limite
all'allungamento o accorciamento di queste fasi medianto il parametro
di \textit{Syncronization Jump Width} \textit{SJW} che può variare da
1 a 4 TQ; questo procedimento è detto \textit{bit resynchronization}.
Un metodo alternativo è quello di resettare il bit time ogni volta
che si presenza un fronte, tecnica denominata \textit{hard
synchronization}
\end{description}

%FIXME immagine per spiegare meglio

\subsection{Sottolivello PMA}
La sigla PMA corrisponde a \textit{Physical Medium Attachment}. Lo
scopo di questo livello è definire come i nodi vengono collegati al
bus. Il PMA prevede 3 componenti:
\begin{description}
 \item[Transceiver] è un buffer bidirezionale per l'accesso al bus
quindi rileva la differenza fra CAN\_H e CAN\_L
 \item[Controller] è il componente che si occupa della trasmissione e
ricezione dei dati sul bus
 \item[Microcontroller] è il dispositivo che si occupa di gestire
transceiver e controller, quindi di decidere cosa trasmettere e
analizzare i dati ricevuti.
\end{description}

%FIXME immagine

\subsection{Sottolivello MDI}
La sigla MDI corrisponde a \textit{Medium Dependent Interface}. A
questo livello sono definite le caratteristiche elettriche e fisiche:
come sono connessi i nodi, impedenze di adattamento, ritardi di linea
ecc.

\section{Livello Data-link}
Il livello data-link si occupa della definizione di come sono
strutturati i messaggi, di come vengono trattati gli errori, di come
vengono convalidati i messaggi, di come isolare i nodi malfunzionanti
e delle regole di arbitraggio.
\subsection{Formato messaggi}
Il protocollo CAN riconosce 5 tipi di formati messaggi:
\begin{description}
 \item[Data Frame] è un messaggio che contiene dati utili; ogni
ricevente può decidere se usarli o scartarli.
 \item[Remote Frame] è un messaggio simile al data frame, ma privo
della parte dati; esso è utilizzato come \textit{data request}.
 \item[Error Frame] è un messaggio inviato da un nodo che rileva un
errore e provoca il re-invio di un messaggio. Se un nodo rileva un
alto tasso di errore si esclude dalla rete.
 \item[Overload] è un messaggio inviato da un nodo occupato per
chiedere di ritardare la spedizione del prossimo messaggio
 \item[Interframe Space] è un messaggio che precede ciascun data
frame e remote frame di modo da permettere la separazione fra questi
due tipi di messaggio.
\end{description}

\subsubsection{Data Frame e Remote Frame}
Un \textit{Data Frame} o \textit{Remote Frame} è costituito da 10
campi:

\begin{description}
 \item[SOF]: Start Of Frame, 1 bit dominante per per indicare
l'inizio del frame
 \item[Arbitration field]: può essere lungo 11 o 29 bit; esso è un
identificatore di messaggio
 \item[RTR]: Remote Transmission Request, 1bit che serve a distinguere
fra Data Frame e Remote Frame; è dominante nel caso del Data Frame,
mentre è recessivo nel caso del Remote Frame. In questo modo la
trasmissione dati ha priorità maggiore della richiesta.
 \item[IDE]: 1 bit riservato
 \item[r0]: 1 bit riservato
 \item[DLC]: 4 bit per indicare la lunghezza del campo Data Field
 \item[Data Field]: da 0 a 8 Byte di dati. Nel caso del Remote Frame,
questo campo è assente, quindi DLC è 0.
 \item[CRC]: 15bit di CRC più uno aggiuntivo sempre recessivo.
 \item[ACK]: 2 bit recessivi, denominati \textit{Slot} e
\textit{Delimiter}; se almeno un nodo destinatario ha ricevuto il
messaggio \textit{Slot} viene da esso sovrascritto con un bit
dominante. Questo garantisce che il messaggio è stato consegnato,
anche se non sappiamo chi ha ricevuto il messaggio.
 \item[EOF]: End Of Frame, 7 bit recessivi indicano la fine del frame.
\end{description}

\subsubsection{Interframe Space}
Come accennato è frame usato per separare i Data Frame dai Remote
Frame. Questo frame è costituito da una sequenza di bit tutti
recessivi. Durante la trasmissione dei primi tre bit recessivi, detti
\textit{intermission}, nessun nodo può inviare Data Frame o Remote
Frame; tuttavia possono essere inviati frame di altro tipo. A seguire
una sequenza di 8 bit recessivi utilizzati da un nodo soggetto ad
errori come attesa per inviare nuovamente il proprio messaggi; altri
nodi però possono iniziare una nuova trasmissione, quindi il nodo
soggetto ad errore può decidere se accettare il nuovo frame in
trasmissione o rifiutarlo a secondo che il suo stato sia,
rispettivamente, \textit{error passive} o \textit{non error passive}.
Infine una sequenza di lunghezza non specificata di bit recessivi per
indicare lo stato di \textit{Bus Idle}, ovvero non ci sono
trasmissioni in corso.

\subsubsection{Overload}
Questo tipo di messaggio consiste in una sequenza di bit il cui scopo
è ritardare l'emissione di Data Frame o Remote Frame; questo consente
ai nodi in stato \textit{busy} di terminare la propria elaborazione e
quindi poter accettare nuovi messaggi.

Un Overload Frame può essere inviato solo durante i primi 3 bit di un
Interframe.

Il messaggio è costituito da due campi. Il primo campo è detto
\textit{flag} e consiste in 6bit dominanti che possono essere inviati
durante la trasmissione dell'intermission di un Interframe. Tutti i
nodi rilevano la situazione di overload e rispondono con un flag di
overload. Se il flag viene inviato dopo l'intermission, questi
vengono interpretati come parte di un Data Frame o Remote Frame, ma
genereranno degli errori.

Il secondo campo è detto \textit{delimiter}; esso consiste in 8 bit
recessivi che indicano la fine del frame.

Un frame di overload ritarda la comunicazione di tutti i nodi
pertanto non è permesso avere più di due overload consecutivi da
parte di uno stesso nodo, il quale altrimenti potrebbe generare uno
stallo su tutto il bus.

\subsubsection{Rilevamento degli errori e Error Frame}
Le cause di errore possono essere di diverso tipo: sfasamento del
sample-point, soglie di riconoscimento differenti, attenuazione del
segnale. Quando un controller rileva un errore invia un Error Frame,
quindi tutti i nodi collegati ignorano il messaggio non valido;
quindi il controller ritrasmette il messaggio originale.

Un controller rileva errori in diverse situazioni:
\begin{description}
 \item[Bit monitoring].Quando un nodo vince l'arbitraggio e i suoi bit
inviati sono diversi da quelli che effettivamente si trovano
sul bus, si ha un \textit{bit error}.
 \item[Bit stuffing]. Se sul bus vengono rilevati più di 6 bit
consecutivi uguali, si ha uno \textit{stuff error}. Come spiegato in
precedenza, ogni 5bit ne viene inviato uno di polarità opposta, se
questo non accade si è in una situazione di errore.
 \item[Frame check]. Si ha quando il formato non corrisponde ad uno
dei tipi di messaggio riconosciuti.
 \item[CRC]. Il codice di CRC ricevuto non corrisponde alla sequenza
di bit ricevuti.
 \item[Acknowledgment]. Se nessun nodo imposta il bit di ack a
dominante, significa che nessuno ha ricevuto il messaggio, quindi si
rileva una situazione di errore.
\end{description}

Ogni nodo tiene due contatori di errori: uno per contare gli errori
di ricezione, l'altro per quelli in trasmissione. Ogni volta che
viene rilevato un errore in ricezione il corrispondente contatore
viene incrementato di un'unità; se l'errore avviene in fase di
trasmissione il corrispondente contatore viene incrementato di 8. In
caso di comunicazione avvenuta con successo il rispettivo contatore
viene decrementato di un'unità.

Fin tanto che entrambe i contatori sono inferiori a 128, il nodo si
pone nello stato di \textit{error active}, quindi il nodo può inviare
flag di errore. Se uno dei due contatori supera la soglia dei 128,
allora il nodo è soggetto a frequenti errori quindi entra nello stato
\textit{error passive} quindi non può effettuare una ritrasmissione
del messaggio soggetto a errore, inoltre dovrà attendere un tempo
maggiore per trasmettere. Infine, se un contatore raggiunge il valore
di 256, allora il nodo è da considerarsi escluso dal bus perché
soggetto a troppi errori, quindi si trova nello stato \textit{bus
off}. Quando un nodo è nello stato di bus off è comunque possibile
riabilitarlo attraverso opportuni comandi da parte dell'utente.

Quando un nodo rileva un errore abbiamo detto che invia un Error
Frame; questo frame è costituito da due campi come l'Overload frame:
flag e delimeter. Il campo flag consiste in 6bit dominanti o
recessivi. Se il nodo è nello stato di active error, userà bit
dominanti di modo da sovrascrivere il bus e generare situazioni di
errore anche negli altri nodi a causa del bit stuffing violato. Se
invece si trova nello stato passive error, i bit inviati saranno di
tipo recessivo, quindi non sovrascrivono il traffico.

Nel caso di active error, quando gli altri nodi rilevano un errore,
invieranno a loro volta un Error Frame, questo porta ad avere un
massimo di 12bit dominanti consecutivi. Dopo questa sequenza di bit
di flag, abbiamo 8bit di delimiter che indicano la fine del messaggio
di errore.

\subsection{Arbitraggio}
Durante la trasmissione di Data Frame o Remote Frame, subito dopo il
bit di SOF, viene inviata una sequenza di bit che indicano
l'identificativo. Come già detto questa sequenza di bit è univoca sul
bus, quindi due nodi non possono inviare lo stesso identificativo.
Durante la trasmissione dell'identificativo viene deciso chi
trasmette e chi ascolta. Se più nodi trasmettono contemporaneamente,
quando uno dei due nodi trasmette un dominante e l'altro trasmette
un recessivo, il primo nodo vince l'arbitraggio e il secondo rimane
in ascolto sul bus.

\end{document}
